\rok{Март 2026.}{А}

Трећи поправни тест из Алгоритама и структура података

Први практични тест одржан 06.03.2026.

\zadatak{1}{Quick Sort}
Написати \texttt{quickSort} методу Quick Sort алгоритма за сортирање целобројног низа дужине $n$.

\textbf{Додатне напомене за решавање задатка:}
\begin{itemize}
	\item Улаз је несортиран низ целих бројева.
	\item Излаз је сортирани низ целих бројева.
	\item У template-у се налази класа \texttt{RandomNumberGen} коју треба искористити за генерисање низа случајних бројева.
	\item У оквиру класе \texttt{QuickSort} написати имплементацију за методу:
	\begin{Verbatim}[fontsize=\small]
private static void quickSort(int[] arr, int lower, int upper)
	\end{Verbatim}
	\item Обезбедити код у Main методи за тестирање алгоритма.
\end{itemize}

\zadatak{2}{Избацивање елемената ван парних граница}
Имплементирати методу која у једном пролазу кроз повезану листу проналази први и последњи парни број. Метода треба да обрише све елементе који се налазе пре првог парног броја и све елементе који се налазе после последњег парног броја. Ако у листи нема парних бројева, листа треба да постане празна.

\textbf{Додатни захтеви за решавање задатка:}
\begin{itemize}
	\item Ако минимална вредност не постоји (празна листа), вратити 0.
	\item Алгоритам писати у оквиру класе \texttt{LinkedList}.
	\item Алгоритам мора имати временску комплексност $O(n)$.
	\item Захтеван потпис методе:
	\begin{Verbatim}[fontsize=\small]
public void trimToEvenBounds()
	\end{Verbatim}
\end{itemize}

\textbf{Пример:}
\begin{itemize}
	\item \textbf{Улаз:} \texttt{1 -> 3 -> 4 -> 5 -> 6 -> 7}
	\item \textbf{Излаз:} \texttt{4 -> 5 -> 6}
\end{itemize}

\zadatak{3}{Monitoring система за кеширање података (Cache Loss Tracker)}
Развијате систем за кеширање који складишти податке различитих величина (у мегабајтима). Кеш меморија има строго ограничен капацитет. Подаци се уписују редом, а систем функционише по принципу "најстарији први напоље" (FIFO).

Опис задатка: Написати методу која прима низ целих бројева \texttt{dataSizes} (величине података који пристижу) и \texttt{cacheCapacity} (максимални капацитет кеша).

\begin{enumerate}
	\item Подаци се додају у кеш један по један.
	\item Пре него што се податак дода, мора се ослободити довољно простора брисањем најстаријих података.
	\item Ако је сам податак који пристиже већи од укупног капацитета кеша, он се одбацује и сматра се "изгубљеним".
	\item Метода треба да израчуна и врати укупну количину података (у MB) који су трајно обрисани или одбачени током процеса.
\end{enumerate}

\textbf{Додатни захтеви за решавање задатка:}
\begin{itemize}
	\item Повратна вредност: \texttt{int}.
	\item Користити помоћну структуру реда (Queue) која је дата у шаблону.
	\item Алгоритам писати у оквиру методе:
	\begin{Verbatim}[fontsize=\small]
public int calculateTotalDataLoss(int[] dataSizes, int cacheCapacity)
	\end{Verbatim}
\end{itemize}

\textbf{Пример:}
\begin{itemize}
	\item \textbf{Улаз 1:} \texttt{dataSizes = [10, 20, 30, 5, 45], cacheCapacity = 50}
	\item \textbf{Симулација корак по корак:}
	\begin{enumerate}
		\item Стиже \texttt{10}: Кеш \texttt{(10/50)}. Loss = \texttt{0}.
		\item Стиже \texttt{20}: Кеш \texttt{(10, 20 / Sum: 30/50)}. Loss = \texttt{0}.
		\item Стиже \texttt{30}: Премашује \texttt{50} (\texttt{30+30=60}). Избацује се најстарији \texttt{10}. Кеш \texttt{(20, 30 / Sum: 50/50)}. Loss = \texttt{10}.
		\item Стиже \texttt{5}: Премашује \texttt{50} (\texttt{50+5=55}). Избацује се најстарији \texttt{20}. Кеш \texttt{(30, 5 / Sum: 35/50)}. Loss = \texttt{10 + 20 = 30}.
		\item Стиже \texttt{45}: Премашује \texttt{50} (\texttt{35+45=80}). Избацује се најстарији \texttt{30}, па следећи \texttt{5}. Кеш \texttt{(45 / Sum: 45/50)}. Loss = \texttt{30 + 30 + 5 = 65}.
	\end{enumerate}
	\item \textbf{Излаз 1:} \texttt{65}
\end{itemize}

\sablon{Шаблон трећег поправног теста март 2026. група А}{2025/Mart/GrupaA/AISP2025_PopravniTest3_Test1_GrupaA.linq}
